\chapter{Invariancia gauge y necesidad del potencial vectorial}

Antes de empezar con la física, hay dos cuestiones prácticas. Se suspenderá la próxima sesión mientras el profesor esté en Bélgica, pero el curso continuará hasta que se hayan impartido las diez clases. Además, los apuntes manuscritos de esta sesión se han preparado con poca antelación; quien quiera escanearlos para la página web de la clase los rescatará de un escáner extraordinariamente lento.

La clase anterior fue densa en matemáticas. En vez de fingir lo contrario, volveremos a desarrollar su cálculo central. El objetivo del curso es ofrecer un verdadero mínimo teórico: la estructura suficiente para avanzar al siguiente nivel. Si se pudiera eliminar la matemática sin perder la física, habría motivos de sobra para hacerlo. Pero las leyes exactas exigen enunciados cuantitativos, y estos requieren matemáticas. La notación puede volverse arcana; eso no demuestra que los físicos intenten ocultar la materia tras un lenguaje sacerdotal.

Quizá aún esté por descubrirse una formulación más sencilla de la naturaleza. Como bromeaba Feynman, se puede elevar al cuadrado cada ecuación de la física, sumar los cuadrados e igualar la suma a cero. Así se obtiene una sola ecuación, porque entonces todos los cuadrados deben anularse, pero no se obtiene una comprensión nueva. La sencillez consiste en encontrar la idea organizadora, no solo en comprimir la tipografía. Cuando alguien del público objeta que esa idea necesariamente debe venir de una persona joven, se concede enseguida que la edad es irrelevante para la intuición.

\lectureoverview{Primero recuperaremos la notación de índices necesaria para reconocer los escalares de Lorentz. Después formularemos cuatro principios para una teoría relativista, repetiremos íntegramente la derivación de la partícula cargada y, por último, descubriremos dos veces la invariancia gauge: como término de extremo en la acción y como invariancia del tensor del campo electromagnético.}

\section{Repaso del lenguaje de índices}

\subsection{Por qué algunos índices van arriba}

Sea
\begin{equation}
x^\mu=(t,x,y,z),
\qquad
\mu=0,1,2,3,
\end{equation}
con métrica
\begin{equation}
\eta_{\mu\nu}
=
\operatorname{diag}(-1,1,1,1).
\label{eq:sr7-metric}
\end{equation}
Las coordenadas y los desplazamientos de coordenadas llevan índices superiores:
\begin{equation}
dx^\mu=(dt,dx,dy,dz).
\end{equation}
Las derivadas respecto de esas coordenadas llevan índices inferiores:
\begin{equation}
\partial_\mu
\equiv
\frac{\partial}{\partial x^\mu}.
\label{eq:sr7-gradient}
\end{equation}
Qué familia se denomina superior y cuál inferior es una convención, como lo son los nombres izquierda y derecha. Sin embargo, una vez elegida, hay que seguirla con coherencia, porque ambas familias se transforman de manera distinta.

\begin{classroomqa}
\audiencequestion{¿Cuándo se escribe un índice como superíndice y cuándo como subíndice?}
\lecturerresponse{A las coordenadas, los desplazamientos y los objetos que se transforman como ellos se les asignan índices superiores. A las derivadas y a los objetos que se transforman como ellas se les asignan índices inferiores. La elección de cuál familia se llama superior es convencional, pero la diferencia entre las dos leyes de transformación no lo es. \lecturetimestamp{00:05:03}}
\end{classroomqa}

Para una función escalar \(s(x)\), el diferencial total es
\begin{equation}
ds
=
\frac{\partial s}{\partial x^\mu}\,dx^\mu
=
\partial_\mu s\,dx^\mu.
\label{eq:sr7-total-differential}
\end{equation}
Escrito término a término, queda
\begin{equation}
ds
=
\frac{\partial s}{\partial t}\,dt
+\frac{\partial s}{\partial x}\,dx
+\frac{\partial s}{\partial y}\,dy
+\frac{\partial s}{\partial z}\,dz.
\end{equation}
La expresión compacta dice exactamente lo mismo. No es una ley nueva, sino un recurso para ver la estructura escalar sin arrastrar cuatro términos en cada línea.

\begin{figure}[htbp]
\centering
\includegraphics[width=0.90\textwidth]{lecture_07_figure_01.png}
\caption{El diferencial escalar formado a partir de un gradiente con índice inferior y un desplazamiento con índice superior.}
\lectureframe{00:08:20}
\end{figure}

La convención de Einstein se entiende ahora de inmediato: se suma sobre todo índice que aparece una vez arriba y otra abajo. Por tanto, la contracción de
\eqref{eq:sr7-total-differential} es un escalar. Los vectores y los covectores tienen una interpretación geométrica más profunda, pero dejaremos ese tema para relatividad general.

\subsection{Subir un índice}

Todo cuadrivector tiene formas con índices superiores e inferiores, relacionadas por la métrica:
\begin{equation}
A_\mu=\eta_{\mu\nu}A^\nu.
\end{equation}
Con \eqref{eq:sr7-metric},
\begin{equation}
A_0=-A^0,
\qquad
A_m=A^m,
\qquad m=1,2,3.
\label{eq:sr7-lowering}
\end{equation}
Así, subir o bajar un índice cambia el signo de la componente temporal y no afecta a las componentes espaciales. La misma regla se aplica a la derivada:
\begin{equation}
\begin{aligned}
\partial_\mu
&=
\left(
\frac{\partial}{\partial t},
\frac{\partial}{\partial x},
\frac{\partial}{\partial y},
\frac{\partial}{\partial z}
\right),\\[0.5em]
\partial^\mu
&=
\left(
-\frac{\partial}{\partial t},
\frac{\partial}{\partial x},
\frac{\partial}{\partial y},
\frac{\partial}{\partial z}
\right).
\end{aligned}
\label{eq:sr7-raised-gradient}
\end{equation}

\begin{figure}[htbp]
\centering
\includegraphics[width=0.90\textwidth]{lecture_07_figure_02.png}
\caption{Los operadores de derivada con índices inferior y superior; solo cambia el signo de la componente temporal.}
\lectureframe{00:12:30}
\end{figure}

\begin{classroomqa}
\audiencequestion{¿Una derivada parcial lleva siempre un índice inferior, y también se puede subir?}
\lecturerresponse{El operador \(\partial/\partial x^\mu\) es naturalmente
\(\partial_\mu\), un covector. Al aplicar la métrica se obtiene
\(\partial^\mu=\eta^{\mu\nu}\partial_\nu\); en el espacio-tiempo plano esto solo invierte el signo de la componente temporal. \lecturetimestamp{00:11:00}}
\end{classroomqa}

El intervalo invariante es el ejemplo conocido:
\begin{equation}
x_\mu x^\mu
=
x^2+y^2+z^2-t^2.
\label{eq:sr7-interval}
\end{equation}
La contracción de un índice superior con uno inferior produce un escalar de Lorentz. Esta sencilla regla de reconocimiento es la razón por la que hay que llevar la notación con tanto cuidado.

\section{Cuatro principios para construir una teoría}

Supongamos que un experimento revela un fenómeno nuevo y queremos formular una teoría que lo describa. Cuatro principios organizan la búsqueda. Quizá haya que refinarlos con el tiempo, pero por ahora todas las teorías fundamentales que han tenido éxito siguen este patrón.

\subsection{El principio de acción}

La primera exigencia es que haya una acción. Para la mecánica de partículas,
\begin{equation}
I[x]
=
\int dt\,L(x,\dot x),
\label{eq:sr7-particle-action-general}
\end{equation}
y la trayectoria física es estacionaria bajo variaciones que mantienen fijos sus extremos:
\begin{equation}
\delta I=0.
\end{equation}
Las ecuaciones resultantes son
\begin{equation}
\frac{d}{dt}
\left(
\frac{\partial L}{\partial\dot x^m}
\right)
=
\frac{\partial L}{\partial x^m}.
\label{eq:sr7-EL}
\end{equation}

El principio de acción no es un adorno estético. Es el marco que relaciona las simetrías con las leyes de conservación. La simetría de una acción bajo traslaciones temporales implica la conservación de la energía; las traslaciones espaciales implican la conservación del momento; otras simetrías relacionadas sustentan otras cantidades conservadas. La acción no es por sí sola la ley de conservación, pero es la estructura mediante la cual la simetría se convierte en una.

\subsection{Localidad}

La segunda exigencia es la localidad. Si se perturba un punto de una cuerda de violín, primero responde el segmento vecino. Después responde el siguiente y la perturbación se propaga. El efecto puede llegar finalmente a toda la cuerda, pero la influencia directa a tiempos cortos es local.

Para una partícula, la localidad se codifica mediante un lagrangiano evaluado a lo largo de la trayectoria. Depende de la posición y la velocidad en un mismo instante, no de pares arbitrarios de instantes alejados. Para un campo \(\phi(x)\),
\begin{equation}
I[\phi]
=
\int d^4x\,
\mathcal L\!\left(\phi,\partial_\mu\phi\right).
\label{eq:sr7-local-field-action}
\end{equation}
La densidad en un suceso depende del campo y de sus primeras derivadas allí. No contiene un producto directo del campo en este punto por el campo en otro punto distante y no relacionado.

\begin{figure}[htbp]
\centering
\includegraphics[width=0.90\textwidth]{lecture_07_figure_03.png}
\caption{Acciones de partículas y campos mostradas bajo el principio de localidad.}
\lectureframe{00:19:30}
\end{figure}

Por eso las ecuaciones de movimiento locales son ecuaciones diferenciales. Las derivadas comparan valores infinitesimalmente próximos. Una alternativa no local podría relacionar la aceleración de ahora con una fuerza en un instante completamente distinto, o relacionar velocidades en dos instantes remotos con una fuerza en un tercero. Tales leyes son matemáticamente concebibles, pero la localidad las excluye.

\begin{classroomqa}
\audiencequestion{¿Contradice el entrelazamiento cuántico este principio de localidad?}
\lecturerresponse{No en el sentido dinámico que estamos usando aquí. El entrelazamiento genera correlaciones, pero una operación local de Alicia no puede utilizarse para alterar el estado reducido de Bob ni para enviarle instantáneamente una señal controlable. La analogía de las monedas permite aislar la diferencia: ver una moneda revela cuál es la moneda correlacionada que tiene la otra persona, pero no actúa sobre esa moneda distante. Las correlaciones cuánticas son más ricas que las de esa analogía clásica, pero aun así no permiten señalización superlumínica. \lecturetimestamp{00:20:02}}
\end{classroomqa}

\subsection{Invariancia de Lorentz}

La tercera exigencia es que las leyes tengan la misma forma en todos los sistemas inerciales. La regla práctica de construcción es sencilla: el lagrangiano debe ser un escalar de Lorentz, de modo que la acción sea invariante. Las contracciones entre índices superiores--inferiores, como \(\partial_\mu\phi\,\partial^\mu\phi\), proporcionan los elementos básicos.

En relatividad general, el enunciado se amplía. El integrando es una densidad escalar y la invariancia se extiende de las transformaciones de Lorentz a transformaciones de coordenadas arbitrarias. Este refinamiento no sustituye el principio actual; es su formulación para el espacio-tiempo curvo.

\begin{classroomqa}
\audiencequestion{¿La invariancia de Lorentz es solo una regla para la relatividad especial y no un principio verdaderamente general?}
\lecturerresponse{Para un problema sin gravitación en el espacio-tiempo plano, la invariancia de Lorentz es el enunciado correcto. En el espacio-tiempo curvo, se sustituye el lagrangiano escalar por la densidad escalar apropiada y se exige invariancia bajo transformaciones de coordenadas arbitrarias. Salvo por ese refinamiento, se mantiene la misma filosofía de construcción. \lecturetimestamp{00:24:16}}
\end{classroomqa}

\begin{classroomqa}
\audiencequestion{¿Qué dimensiones tienen el lagrangiano de una partícula y la densidad lagrangiana de un campo?}
\lecturerresponse{El lagrangiano de una partícula tiene dimensiones de energía. En
\(I=\int dt\,d^3x\,\mathcal L\), la medida espacial aporta un volumen, así que
\(\mathcal L\) tiene dimensiones de energía por unidad de volumen. Al integrarla en el espacio, se obtiene un lagrangiano ordinario con dimensiones de energía.
\lecturetimestamp{00:25:49}}
\end{classroomqa}

\subsection{Invariancia gauge}

La cuarta exigencia es la invariancia gauge. Por ahora, se deja deliberadamente sin definir. Alguien del público propone una anticipación útil: quizá se puedan añadir ciertas cosas a un lagrangiano sin cambiar la física. La idea se aproxima, pero la palabra «ciertas» hará casi todo el trabajo.

\begin{figure}[htbp]
\centering
\includegraphics[width=0.90\textwidth]{lecture_07_figure_04.png}
\caption{Localidad, construcción escalar, invariancia de Lorentz e invariancia gauge, reunidas en la pizarra.}
\lectureframe{00:27:00}
\end{figure}

Aplazamos la definición porque el cálculo de la partícula cargada ya la contiene. Repetir ese cálculo no es un desvío: mostrará por qué el potencial vectorial es inevitable y por qué, una vez introducido, contiene una redundancia.

\section{Por qué es necesario el potencial vectorial}

Para una partícula cargada que se mueve lentamente, la ecuación buscada es la conocida ley de la fuerza de Lorentz
\begin{equation}
m\mathbf a
=
e\left(
\mathbf E+\mathbf v\times\mathbf B
\right).
\label{eq:sr7-lorentz-target}
\end{equation}

\begin{figure}[htbp]
\centering
\includegraphics[width=0.90\textwidth]{lecture_07_figure_05.png}
\caption{La conocida ley de la fuerza, escrita antes de volver a introducir el potencial vectorial.}
\lectureframe{00:29:08.870}
\end{figure}

En \eqref{eq:sr7-lorentz-target} no aparece ningún potencial vectorial. Se pueden escribir las ecuaciones de movimiento directamente con \(\mathbf E\) y \(\mathbf B\). La dificultad es más estructural: bajo los supuestos de localidad e invariancia de Lorentz empleados en la clase, la interacción de una partícula cargada no se puede escribir como una acción sobre la línea de universo directamente en términos de esos campos, la posición y la velocidad por sí solos. El potencial \(A_\mu\) es necesario para construir el acoplamiento escalar local.

\begin{classroomqa}
\audiencequestion{Si la ley de la fuerza ya contiene solo \(\mathbf E\) y \(\mathbf B\), ¿por qué introducir \(A_\mu\)?}
\lecturerresponse{Porque una ecuación de movimiento no es todavía un principio de acción. En el marco local e invariante de Lorentz de esta clase, la acción de una partícula cargada requiere el cuadripotencial. Los campos bastan para enunciar la fuerza, pero el potencial es el objeto que permite derivarla de una acción escalar sobre la línea de universo. \lecturetimestamp{00:29:06}}
\end{classroomqa}

\subsection{La acción sobre la línea de universo}

La acción candidata consta de una parte libre y otra de interacción:
\begin{equation}
I[x]
=
-m\int d\tau
-e\int A_\mu(x)\,dx^\mu.
\label{eq:sr7-worldline-action}
\end{equation}
El primer término suma la longitud de tiempo propio de cada segmento infinitesimal de la trayectoria. Es escalar por construcción. En el segundo término, \(A_\mu(x)\) se evalúa en el suceso ocupado por la partícula y se contrae con el desplazamiento \(dx^\mu\). Ese incremento también es escalar y local a lo largo de la línea de universo.

\begin{figure}[htbp]
\centering
\includegraphics[width=0.90\textwidth]{lecture_07_figure_06.png}
\caption{Las dos partes de la acción de la partícula cargada y la contracción local \(A_\mu(x)dx^\mu\).}
\lectureframe{00:31:20}
\end{figure}

\subsection{Reescritura de la acción en tiempo ordinario}

Para usar las conocidas ecuaciones de Euler--Lagrange, dividimos cada incremento de la línea de universo por \(dt\):
\begin{equation}
\frac{d\tau}{dt}
=
\sqrt{1-\dot{\mathbf x}^{\,2}},
\qquad
\dot{\mathbf x}^{\,2}
=
\dot x^2+\dot y^2+\dot z^2.
\end{equation}
Las cuatro componentes de \(dx^\mu/dt\) son
\begin{equation}
\frac{dx^\mu}{dt}
=
\left(1,\dot x,\dot y,\dot z\right).
\end{equation}
Por consiguiente,
\begin{align}
I
&=
\int dt\,
\left[
-m\sqrt{1-\dot{\mathbf x}^{\,2}}
-eA_\mu\frac{dx^\mu}{dt}
\right]
\notag\\
&=
\int dt\,
\left[
-m\sqrt{1-\dot{\mathbf x}^{\,2}}
-e\left(A_0+A_m\dot x^m\right)
\right].
\label{eq:sr7-action-in-t}
\end{align}
Así pues, el lagrangiano en tiempo ordinario es
\begin{equation}
\begin{aligned}
L(x,\dot x,t)
&=
-m\sqrt{1-\dot{\mathbf x}^{\,2}}
\\
&\quad -e\left[
A_0(x,t)+A_m(x,t)\dot x^m
\right].
\end{aligned}
\label{eq:sr7-lagrangian}
\end{equation}

\begin{figure}[htbp]
\centering
\includegraphics[width=0.90\textwidth]{lecture_07_figure_07.png}
\caption{La acción sobre la línea de universo, reescrita como una integral respecto del tiempo ordinario, y el lagrangiano obtenido de su integrando.}
\lectureframe{00:35:30}
\end{figure}

A estas alturas, la parte imaginativa ya ha terminado. Una vez elegido el lagrangiano, lo demás es mecánico. Los cálculos pueden ser largos y tediosos, pero esa labor no es inútil: es lo que revela las estructuras eléctrica, magnética y relativista. El profesor bromea diciendo que descubrir el principio es la parte entretenida, y que deducir todas sus consecuencias es lo que se deja a los estudiantes de posgrado.

\begin{classroomqa}
\audiencequestion{¿Se puede subir o bajar el índice espacial de \(A_m\) o \(\dot x^m\) en este cálculo?}
\lecturerresponse{Sí. La métrica espacial es la identidad, así que subir o bajar un índice latino no cambia el valor de su componente. Hay que tener cuidado con el signo del índice temporal. \lecturetimestamp{00:34:10}}
\end{classroomqa}

\section{El cálculo de Euler--Lagrange}

\subsection{Momento canónico y una corrección en directo}

Elijamos una componente espacial \(x^m\). Al derivar la parte libre de \eqref{eq:sr7-lagrangian} se obtiene
\begin{equation}
\frac{m\dot x_m}{\sqrt{1-\dot{\mathbf x}^{\,2}}}
=
m u_m,
\qquad
u^\mu\equiv\frac{dx^\mu}{d\tau}.
\end{equation}
El primer intento en la pizarra se detiene ahí. Una corrección del público detecta la omisión: el término de interacción también depende de \(\dot x^m\). La derivada completa es
\begin{equation}
\frac{\partial L}{\partial\dot x^m}
=
\frac{m\dot x_m}{\sqrt{1-\dot{\mathbf x}^{\,2}}}
-eA_m
=
m u_m-eA_m.
\label{eq:sr7-canonical-momentum}
\end{equation}

\begin{classroomqa}
\audiencequestion{¿A la primera expresión de \(\partial L/\partial\dot x^m\) le falta la contribución del término electromagnético?}
\lecturerresponse{Sí. Al derivar
\(-eA_n\dot x^n\) respecto de \(\dot x^m\) se obtiene la contribución
\(-eA_m\). El momento canónico corregido es
\(m u_m-eA_m\), y no solo el momento mecánico.
\lecturetimestamp{00:37:24}}
\end{classroomqa}

\begin{figure}[htbp]
\centering
\includegraphics[width=0.90\textwidth]{lecture_07_figure_08.png}
\caption{El momento canónico corregido, con los términos mecánico y del potencial.}
\lectureframe{00:38:30}
\end{figure}

\subsection{Derivación a lo largo de la trayectoria en movimiento}

El lado izquierdo de la ecuación de Euler--Lagrange es
\begin{equation}
\frac{d}{dt}
\left(
\frac{\partial L}{\partial\dot x^m}
\right)
=
m\frac{du_m}{dt}
-e\frac{dA_m(x(t),t)}{dt}.
\end{equation}
El potencial cambia por dos razones: puede depender explícitamente del tiempo y la partícula se desplaza por una región donde varía espacialmente. La regla de la cadena da
\begin{equation}
\frac{dA_m}{dt}
=
\partial_0A_m
+\dot x^n\partial_nA_m.
\label{eq:sr7-chain-rule}
\end{equation}
Escrito sin índices, el segundo término significa
\begin{equation}
\dot x\,\frac{\partial A_m}{\partial x}
+\dot y\,\frac{\partial A_m}{\partial y}
+\dot z\,\frac{\partial A_m}{\partial z}.
\end{equation}

\begin{figure}[htbp]
\centering
\includegraphics[width=0.90\textwidth]{lecture_07_figure_09.png}
\caption{Las contribuciones explícita y convectiva a la derivada temporal de \(A_m\).}
\lectureframe{00:42:10}
\end{figure}

En el lado derecho,
\begin{equation}
\frac{\partial L}{\partial x^m}
=
-e\partial_mA_0
-e\dot x^n\partial_mA_n.
\label{eq:sr7-spatial-derivative-L}
\end{equation}
El término libre no depende explícitamente de la posición. En \eqref{eq:sr7-spatial-derivative-L}, \(m\) es el índice libre que selecciona una ecuación; \(n\) es un índice mudo repetido y, por tanto, se suma sobre él. Cambiar el nombre de ese índice mudo no altera nada. Por eso el profesor se detiene en las letras antes de combinar los términos.

Igualando ambos lados y reordenando se obtiene
\begin{equation}
m\frac{du_m}{dt}
=
e\left(
\partial_0A_m-\partial_mA_0
\right)
+e\dot x^n
\left(
\partial_nA_m-\partial_mA_n
\right).
\label{eq:sr7-grouped-force}
\end{equation}

\begin{figure}[htbp]
\centering
\includegraphics[width=0.90\textwidth]{lecture_07_figure_10.png}
\caption{La fuerza separada en un corchete independiente de la velocidad y otro antisimétrico, lineal en la velocidad.}
\lectureframe{00:49:30}
\end{figure}

\subsection{Identificación de los campos eléctrico y magnético}

El primer corchete no contiene la velocidad de la partícula. Según la convención de esta acción, definamos
\begin{equation}
E_m
\equiv
\partial_0A_m-\partial_mA_0.
\label{eq:sr7-E-from-A}
\end{equation}
Su función operativa es precisamente la del campo eléctrico. El segundo corchete es antisimétrico en sus índices espaciales:
\begin{equation}
C_{mn}
\equiv
\partial_nA_m-\partial_mA_n.
\label{eq:sr7-C}
\end{equation}
Sus tres componentes independientes son el rotacional del potencial espacial, con la orientación elegida en la convención de la clase de modo que
\begin{equation}
\dot x^n C_{mn}
=
\left(\mathbf v\times\mathbf B\right)_m.
\label{eq:sr7-magnetic-identification}
\end{equation}
Se deja en la pizarra como ejercicio comprobar los componentes en detalle. La identificación estructural ya es evidente: una derivada antisimétrica de un vector tiene tres componentes, exactamente el número de componentes de \(\mathbf B\), y multiplica a la velocidad.

\begin{figure}[htbp]
\centering
\includegraphics[width=0.90\textwidth]{lecture_07_figure_11.png}
\caption{Los dos corchetes de derivadas identificados como las contribuciones eléctrica y magnética a la fuerza de Lorentz.}
\lectureframe{00:52:00}
\end{figure}

Así pues, la ecuación \eqref{eq:sr7-grouped-force} ha recuperado la fuerza de Lorentz a partir de la acción en términos del potencial vectorial. Además, ha revelado algo aún más importante: las derivadas antisimetrizadas que se convertirán en el tensor del campo.

\section{Del tiempo ordinario a una ecuación de cuadrivectores}

\subsection{La conversión al tiempo propio}

El lado izquierdo de \eqref{eq:sr7-grouped-force} no es la aceleración newtoniana \(d\mathbf v/dt\). Es la derivada respecto del tiempo ordinario de una componente de la cuadrivelocidad:
\begin{equation}
\frac{du^m}{dt}
=
\frac{du^m}{d\tau}\frac{d\tau}{dt}.
\end{equation}
Multipliquemos la ecuación por \(dt/d\tau\). Entonces
\begin{equation}
\dot x^n\frac{dt}{d\tau}
=
\frac{dx^n}{dt}\frac{dt}{d\tau}
=
u^n,
\qquad
\frac{dt}{d\tau}=u^0.
\end{equation}
La ecuación espacial queda
\begin{align}
m\frac{du^m}{d\tau}
&=
e\left(
\partial_0A_m-\partial_mA_0
\right)u^0
\notag\\
&\quad
+e\left(
\partial_nA_m-\partial_mA_n
\right)u^n.
\label{eq:sr7-proper-spatial}
\end{align}

\begin{figure}[htbp]
\centering
\includegraphics[width=0.90\textwidth]{lecture_07_figure_12.png}
\caption{Los factores \(dt/d\tau\) convierten las velocidades ordinarias en componentes de la cuadrivelocidad.}
\lectureframe{00:56:30}
\end{figure}

El lado izquierdo es la aceleración propia:
\begin{equation}
\frac{du^\mu}{d\tau}
=
\frac{d^2x^\mu}{d\tau^2}.
\end{equation}
Es un cuadrivector. Por tanto, el lado derecho debe construirse a partir de un tensor de rango dos que actúe sobre \(u^\nu\).

\subsection{El tensor del campo y la ley de la fuerza}

Elijamos la convención
\begin{equation}
F_{\mu\nu}
\equiv
\partial_\mu A_\nu-\partial_\nu A_\mu.
\label{eq:sr7-F-definition}
\end{equation}
Es antisimétrico:
\begin{equation}
F_{\mu\nu}=-F_{\nu\mu}.
\end{equation}
Con el signo de interacción de \eqref{eq:sr7-worldline-action}, la ecuación compacta es
\begin{equation}
\boxed{
m\frac{du^\mu}{d\tau}
=
-eF^\mu{}_\nu u^\nu
}.
\label{eq:sr7-covariant-force}
\end{equation}
Equivalentemente, antes de denominarlo \(F_{\mu\nu}\),
\begin{equation}
m\frac{du^\mu}{d\tau}
=
e\left(
\partial_\nu A^\mu-\partial^\mu A_\nu
\right)u^\nu.
\end{equation}
Las distintas convenciones para la métrica, la carga y el potencial redistribuyen los signos globales. El contenido invariante es la derivada antisimetrizada del potencial y su contracción con la cuadrivelocidad.

\begin{figure}[htbp]
\centering
\includegraphics[width=0.90\textwidth]{lecture_07_figure_13.png}
\caption{La ecuación en tiempo propio, expresada mediante un tensor del campo de rango dos que actúa sobre la cuadrivelocidad.}
\lectureframe{00:58:25}
\end{figure}

\begin{classroomqa}
\audiencequestion{La derivación produjo tres ecuaciones espaciales. ¿Por qué se puede añadir una cuarta ecuación?}
\lecturerresponse{La acción escalar garantiza un resultado covariante de Lorentz. Una vez que se reconoce que la diferencia entre ambos lados es un cuadrivector, las transformaciones de Lorentz mezclan sus componentes espaciales y temporales. Una ecuación de cuadrivectores no puede tener coherentemente solo tres componentes: la componente temporal acompaña a las espaciales. \lecturetimestamp{00:59:40}}
\end{classroomqa}

\begin{classroomqa}
\audiencequestion{¿Una lista fija de números como \((1,2,3,4)\) es por sí misma un cuadrivector invariante de Lorentz?}
\lecturerresponse{No. Esos números pueden ser las componentes de un vector en un sistema de referencia elegido, pero las componentes se transforman. Incluso en dos dimensiones, el vector de componentes \((1,2)\) se puede describir en un sistema de ejes rotado como \((0,\sqrt5)\). El vector geométrico permanece; su lista numérica de componentes no. \lecturetimestamp{01:00:28}}
\end{classroomqa}

\begin{figure}[htbp]
\centering
\includegraphics[width=0.90\textwidth]{lecture_07_figure_14.png}
\caption{Un vector de dos componentes expresado de nuevo en unos ejes alineados con el propio vector.}
\lectureframe{01:02:00}
\end{figure}

\begin{classroomqa}
\audiencequestion{Si el objetivo era sustituir el tiempo ordinario por el tiempo propio, ¿por qué sigue apareciendo \(\partial_0A_m=\partial A_m/\partial t\)?}
\lecturerresponse{El objetivo no era eliminar toda aparición de la letra \(t\). Era escribir la dinámica en una notación manifiestamente cuadrivectorial. El operador \(\partial_0\) es la componente temporal del cuadrigradiente covariante, del mismo modo que \(\partial_m\) son sus componentes espaciales. Forma parte del tensor cuadridimensional. \lecturetimestamp{01:05:10}}
\end{classroomqa}

La clase aprovecha este momento para volver a enunciar la localidad. Una ecuación local relaciona la aceleración en un suceso con el campo y la velocidad en ese mismo suceso. No vincula la aceleración de ahora con una fuerza en un instante remoto.

\begin{figure}[htbp]
\centering
\includegraphics[width=0.90\textwidth]{lecture_07_figure_15.png}
\caption{La ecuación covariante de la fuerza, utilizada para reformular la localidad como dinámica diferencial.}
\lectureframe{01:05:50}
\end{figure}

\subsection{Representación de las componentes electromagnéticas}

Para el tensor con todos los índices superiores, la convención empleada en la pizarra da
\begin{equation}
F^{\mu\nu}
=
\begin{pmatrix}
0&-E_x&-E_y&-E_z\\
E_x&0&B_z&-B_y\\
E_y&-B_z&0&B_x\\
E_z&B_y&-B_x&0
\end{pmatrix}.
\label{eq:sr7-F-matrix}
\end{equation}
El índice de fila es \(\mu\) y el de columna es \(\nu\). Las entradas mixtas de tiempo--espacio contienen el campo eléctrico; el bloque espacio--espacio contiene el campo magnético. La antisimetría hace que la diagonal sea nula y reduce las dieciséis entradas mostradas a seis componentes independientes.

\begin{figure}[htbp]
\centering
\includegraphics[width=0.90\textwidth]{lecture_07_figure_16.png}
\caption{El tensor electromagnético completo, con las componentes eléctricas en las posiciones de tiempo--espacio y las magnéticas en el bloque espacial.}
\lectureframe{01:12:10}
\end{figure}

\begin{classroomqa}
\audiencequestion{¿Existe una regla estándar para decidir qué índice de un tensor señala las filas y cuál las columnas?}
\lecturerresponse{Es una convención de representación. Aquí, el primer índice \(\mu\) se dispone en vertical y el segundo, \(\nu\), en horizontal. Una vez anunciada esta convención, cada componente se puede leer sin ambigüedad. Subir un índice temporal introduce un signo, pero no altera la convención de filas--columnas. \lecturetimestamp{01:09:01}}
\end{classroomqa}

\begin{classroomqa}
\audiencequestion{¿Qué posiciones de los índices conservan la antisimetría?}
\lecturerresponse{Los tensores \(F_{\mu\nu}\) y \(F^{\mu\nu}\) son antisimétricos. La matriz mixta \(F^\mu{}_\nu\) no es antisimétrica bajo la transposición matricial ordinaria: subir un índice temporal cambia las entradas eléctricas y hace simétrico el par tiempo--espacio, mientras que el bloque magnético puramente espacial sigue siendo antisimétrico. El profesor llega a esta conclusión después de probar brevemente en la pizarra la posibilidad contraria y descartarla.
\lecturetimestamp{01:13:18}}
\end{classroomqa}

\section{¿Qué se entiende por invariancia?}

Ahora sí volvemos al cuarto principio. Una invariancia es un cambio que deja inalteradas las ecuaciones de movimiento. La acción puede permanecer exactamente igual o cambiar en un término de borde que no afecta a las ecuaciones variacionales.

\subsection{Desplazamientos constantes de un campo escalar sin masa}

Consideremos la ecuación del campo escalar sin masa en una dimensión espacial,
\begin{equation}
\frac{\partial^2\phi}{\partial t^2}
-\frac{\partial^2\phi}{\partial x^2}
=0,
\label{eq:sr7-wave}
\end{equation}
que se obtiene a partir de
\begin{equation}
\mathcal L_0
=
-\frac12\partial_\mu\phi\,\partial^\mu\phi.
\label{eq:sr7-massless-L}
\end{equation}
Si \(\phi\) es una solución, también lo es
\begin{equation}
\phi\longrightarrow\phi+c,
\label{eq:sr7-constant-shift}
\end{equation}
donde \(c\) es constante. Todas las derivadas de \(c\) se anulan, así que ni la ecuación ni el lagrangiano que depende solo de derivadas cambian.

\begin{figure}[htbp]
\centering
\includegraphics[width=0.90\textwidth]{lecture_07_figure_17.png}
\caption{La ecuación de ondas sin masa, su lagrangiano en términos de derivadas y el desplazamiento constante \(\phi\mapsto\phi+c\).}
\lectureframe{01:17:10}
\end{figure}

Añadamos ahora un término de masa:
\begin{equation}
\frac{\partial^2\phi}{\partial t^2}
-\frac{\partial^2\phi}{\partial x^2}
+\mu^2\phi
=0,
\end{equation}
\begin{equation}
\mathcal L_\mu
=
-\frac12\partial_\mu\phi\,\partial^\mu\phi
-\frac{\mu^2}{2}\phi^2.
\label{eq:sr7-massive-L}
\end{equation}
La parte con derivadas sigue siendo invariante bajo \(\phi\mapsto\phi+c\), pero \(\phi^2\) no. Ambas teorías son legítimas; una tiene la invariancia bajo desplazamientos y la otra no.

\begin{figure}[htbp]
\centering
\includegraphics[width=0.90\textwidth]{lecture_07_figure_18.png}
\caption{El término de masa que rompe la invariancia bajo desplazamientos constantes.}
\lectureframe{01:18:50}
\end{figure}

\subsection{Un estudiante se adelanta al término de borde}

Alguien del público propone una prueba más precisa:
\begin{equation}
\phi(x,t)\longrightarrow\phi(x,t)+x.
\label{eq:sr7-linear-shift}
\end{equation}
Las segundas derivadas de \(x\) se anulan, así que la ecuación de ondas sin masa sigue cumpliéndose. La acción requiere un examen más atento. Puesto que \(x\) y \(t\) son variables independientes,
\begin{equation}
\frac{\partial x}{\partial t}=0,
\qquad
\frac{\partial x}{\partial x}=1.
\end{equation}
La parte espacial cambia de acuerdo con
\begin{align}
-\frac12
\left[
\partial_x(\phi+x)
\right]^2
&=
-\frac12
\left(
\partial_x\phi+1
\right)^2
\notag\\
&=
-\frac12(\partial_x\phi)^2
-\partial_x\phi
-\frac12.
\label{eq:sr7-linear-shift-L}
\end{align}
Los nuevos términos son una constante y una derivada total. La constante no cambia la ecuación de Euler--Lagrange. Al integrar \(\partial_x\phi\) solo se obtienen valores en los extremos espaciales, de modo que, bajo las habituales condiciones de contorno fijas, tampoco cambia la ecuación de movimiento.

\begin{figure}[htbp]
\centering
\includegraphics[width=0.90\textwidth]{lecture_07_figure_19.png}
\caption{El desplazamiento propuesto \(\phi\mapsto\phi+x\) y el cuadrado de la derivada, que revela una constante más una derivada total.}
\lectureframe{01:22:40}
\end{figure}

\begin{classroomqa}
\audiencequestion{¿La transformación \(\phi\mapsto\phi+x\) conserva la ecuación del campo sin masa aunque cambie la densidad lagrangiana?}
\lecturerresponse{Sí. Las segundas derivadas de \(x\) se anulan, así que la ecuación no cambia. El lagrangiano solo cambia en una constante y una derivada espacial total. Estas modifican el valor numérico de la acción en una constante fija de volumen y en valores de borde, pero no la ecuación de Euler--Lagrange bajo condiciones de contorno fijas. \lecturetimestamp{01:20:00}}
\end{classroomqa}

\begin{classroomqa}
\audiencequestion{¿Por qué es cero \(\partial x/\partial t\)?}
\lecturerresponse{Porque \(x\) y \(t\) son coordenadas independientes en la descripción del campo. Una derivada parcial respecto de \(t\) mantiene \(x\) fija. \lecturetimestamp{01:22:22}}
\end{classroomqa}

El estudiante ha llegado al recurso que necesitábamos para el electromagnetismo: una derivada total puede cambiar la forma escrita de la acción sin cambiar sus trayectorias estacionarias.

\section{Invariancia gauge de la partícula cargada}

\subsection{Adición de un cuatrogradiente}

Volvamos al término de interacción
\begin{equation}
I_{\mathrm{int}}
=
-e\int A_\mu\,dx^\mu.
\end{equation}
Sea \(\Lambda(x)\) un escalar suficientemente suave y transformemos el potencial según
\begin{equation}
A_\mu
\longrightarrow
A'_\mu
=
A_\mu+\partial_\mu\Lambda.
\label{eq:sr7-gauge-transformation}
\end{equation}
El gradiente \(\partial_\mu\Lambda\) tiene la misma estructura de índices covariantes que \(A_\mu\), así que la suma está bien definida.

\begin{figure}[htbp]
\centering
\includegraphics[width=0.90\textwidth]{lecture_07_figure_20.png}
\caption{El acoplamiento sobre la línea de universo y la transformación \(A_\mu\mapsto A_\mu+\partial_\mu\Lambda\).}
\lectureframe{01:24:50}
\end{figure}

El cambio en la acción es
\begin{align}
\Delta I_{\mathrm{int}}
&=
-e\int \partial_\mu\Lambda\,dx^\mu
\notag\\
&=
-e\int d\Lambda
\notag\\
&=
-e\left[
\Lambda(x_2)-\Lambda(x_1)
\right].
\label{eq:sr7-gauge-boundary}
\end{align}
A lo largo de cada segmento infinitesimal de la línea de universo,
\(\partial_\mu\Lambda\,dx^\mu\) es exactamente el incremento
\(d\Lambda\). Al sumar los incrementos, la suma telescópica da la diferencia entre los valores en los extremos final e inicial.

\begin{figure}[htbp]
\centering
\includegraphics[width=0.90\textwidth]{lecture_07_figure_21.png}
\caption{La contribución del gradiente, integrada a lo largo del camino, da la diferencia entre los valores en los extremos \(\Lambda(x_2)-\Lambda(x_1)\).}
\lectureframe{01:28:10}
\end{figure}

El principio variacional compara trayectorias cercanas manteniendo fijos \(x_1\) y
\(x_2\). Por tanto, el término adicional en \eqref{eq:sr7-gauge-boundary} es el mismo para toda variación admisible:
\begin{equation}
\delta\Delta I_{\mathrm{int}}=0.
\end{equation}
La trayectoria estacionaria no cambia. Esta es la invariancia gauge entendida como una afirmación sobre los extremos.

\begin{classroomqa}
\audiencequestion{¿Puede \(\Lambda\) ser cualquier función escalar?}
\lecturerresponse{Sí, siempre que se cumplan las condiciones de suavidad y de contorno necesarias para las derivadas y la acción. La transformación es local porque el parámetro es una función arbitraria del espacio-tiempo, y no una única constante global. \lecturetimestamp{01:28:40}}
\end{classroomqa}

Añadir un gradiente a \(A_\mu\) se denomina transformación gauge. La palabra gauge se ha heredado históricamente; en la clase se advierte, con humor pero con firmeza, que no debe pronunciarse como si fuera una palabra francesa. El nombre importa menos que la operación.

\subsection{Demostración directa mediante el tensor de campo}

Las ecuaciones de movimiento dependen de \(F_{\mu\nu}\), no del potencial por sí solo. Bajo \eqref{eq:sr7-gauge-transformation},
\begin{align}
F'_{\mu\nu}
&=
\partial_\mu
\left(
A_\nu+\partial_\nu\Lambda
\right)
-
\partial_\nu
\left(
A_\mu+\partial_\mu\Lambda
\right)
\notag\\
&=
F_{\mu\nu}
+
\partial_\mu\partial_\nu\Lambda
-
\partial_\nu\partial_\mu\Lambda
\notag\\
&=
F_{\mu\nu}.
\label{eq:sr7-F-gauge-invariant}
\end{align}
Para un escalar suave en el espacio-tiempo plano ordinario, las derivadas parciales conmutan. Los dos términos adicionales se cancelan. Así, el tensor de campo, los campos eléctrico y magnético que contiene y la ley de fuerza son todos invariantes gauge.

\begin{figure}[htbp]
\centering
\includegraphics[width=0.90\textwidth]{lecture_07_figure_22.png}
\caption{Las dos segundas derivadas que genera una transformación gauge se cancelan porque las derivadas parciales conmutan.}
\lectureframe{01:33:50}
\end{figure}

Las dos demostraciones completan el mismo argumento desde direcciones opuestas. La acción solo cambia en los extremos fijos, mientras que las ecuaciones dependen únicamente de un tensor de campo que no cambia.

\section{Qué significa la redundancia}

\subsection{¿Por qué conservar un objeto que no es único?}

La objeción natural surge de inmediato. Si \(A_\mu\) no es único, ¿por qué lo introdujimos? Porque es indispensable para la acción local e invariante de Lorentz, aunque los valores de sus componentes contengan información redundante. Muchos potenciales representan el mismo campo electromagnético:
\begin{equation}
A_\mu
\sim
A_\mu+\partial_\mu\Lambda.
\end{equation}
La clase de equivalencia es física; elegir un representante particular es una elección de descripción.

\begin{classroomqa}
\audiencequestion{¿Tiene alguna utilidad práctica cambiar el gauge?}
\lecturerresponse{Sí. Un \(\Lambda\) adecuado puede simplificar una componente del potencial o hacer evidente algún aspecto del problema. Localmente y con condiciones de contorno adecuadas, a menudo se puede elegir un gauge en el que se anule una componente determinada, incluido \(A_0\). Otro gauge puede mostrar con mayor claridad un aspecto diferente. \lecturetimestamp{01:35:43}}
\end{classroomqa}

\begin{classroomqa}
\audiencequestion{¿Es la invariancia gauge análoga a describir el mismo experimento desde dos marcos de Lorentz?}
\lecturerresponse{No. Las transformaciones de Lorentz comparan marcos de referencia físicamente significativos, asociados a distintos observadores. Las transformaciones gauge son redundancias dentro de una misma descripción: distintos potenciales codifican el mismo campo. Matemáticamente, ambas son invariancias, pero sus interpretaciones son distintas. \lecturetimestamp{01:37:07}}
\end{classroomqa}

También hay una diferencia estructural. Una rotación usa un ángulo global fijo; por lo general, no se giran los ejes con ángulos distintos en cada punto. Una transformación gauge se parametriza mediante la función completa \(\Lambda(x)\). Esta libertad dependiente de la posición es lo que la convierte en una invariancia gauge local.

\begin{classroomqa}
\audiencequestion{¿Añadir un gradiente es la definición general de invariancia gauge?}
\lecturerresponse{En la teoría electromagnética abeliana, esa es la transformación. Las teorías de Yang--Mills tienen una versión no abeliana más elaborada, en la que el potencial posee estructura interna adicional, pero la misma idea se mantiene: distintas descripciones locales representan la misma física. \lecturetimestamp{01:39:27}}
\end{classroomqa}

\begin{classroomqa}
\audiencequestion{¿Es una transformación gauge simplemente un recurso matemático para destacar distintas partes de un problema?}
\lecturerresponse{Es un recurso útil para el cálculo y la interpretación, pero su función está limitada: puede simplificar las ecuaciones o revelar una propiedad, pero no puede cambiar ninguna predicción final invariante gauge.
\lecturetimestamp{01:40:22}}
\end{classroomqa}

\subsection{Otros campos y otros potenciales}

Los potenciales gauge de otras teorías también suelen denotarse por \(A\) y llevan un índice de espacio-tiempo \(\mu\). Se distinguen por la estructura interna adicional asociada a otras simetrías. En la clase, esa estructura se deja para más adelante en lugar de presentar una fórmula prematuramente.

\begin{classroomqa}
\audiencequestion{¿Qué aspecto tienen los potenciales vectoriales de otras teorías de campos?}
\lecturerresponse{Siguen comportándose como vectores de espacio-tiempo y suelen escribirse \(A_\mu\), pero llevan etiquetas adicionales y una estructura algebraica asociadas a simetrías internas. Esta es la puerta de entrada a la teoría de Yang--Mills, y se deja intencionadamente para más adelante. \lecturetimestamp{01:41:02}}
\end{classroomqa}

\begin{classroomqa}
\audiencequestion{¿Es realmente imposible escribir directamente con \(\mathbf E\) y \(\mathbf B\) la acción local e invariante de Lorentz de una partícula cargada?}
\lecturerresponse{Dentro de los supuestos aquí indicados, sí. La clase presenta esto como un teorema, pero no lo demuestra. Preservar el principio de acción, la localidad y la invariancia de Lorentz obliga a introducir el potencial; a partir de ahí, la redundancia gauge aparece automáticamente.
\lecturetimestamp{01:42:18}}
\end{classroomqa}

\subsection{¿Qué se puede medir?}

Un valor local de \(A_\mu\) por sí solo no es invariante gauge. Si una medición física determinara ese valor de manera única, cambiarlo en \(\partial_\mu\Lambda\) alteraría un resultado físico, en contradicción con la invariancia que acabamos de demostrar. Las predicciones medibles deben construirse a partir de cantidades invariantes gauge, como \(F_{\mu\nu}\) en esta discusión clásica.

\begin{classroomqa}
\audiencequestion{¿Por qué la naturaleza elige teorías gauge con tanta insistencia?}
\lecturerresponse{En el nivel clásico, la clase no ofrece una razón definitiva. La mecánica cuántica aporta una perspectiva más profunda, pero ni siquiera allí se agota la respuesta general. Lo que ya resulta evidente es la singular belleza, simetría y economía de las teorías organizadas mediante redundancia gauge.
\lecturetimestamp{01:43:36}}
\end{classroomqa}

\begin{classroomqa}
\audiencequestion{¿Se puede medir el propio potencial vectorial?}
\lecturerresponse{No como valor local de una componente dependiente del gauge. Una medición es un proceso físico, y toda predicción física debe ser insensible a \(A_\mu\mapsto A_\mu+\partial_\mu\Lambda\). Lo que puede medirse son propiedades invariantes gauge del potencial y del campo.
\lecturetimestamp{01:44:18}}
\end{classroomqa}

\begin{classroomqa}
\audiencequestion{¿Qué ocurre si \(A_\mu\) cambia en algo que no es un gradiente?}
\lecturerresponse{En general, \(F_{\mu\nu}\) cambia, por lo que el cambio es detectable. En cambio, un potencial de gauge puro \(A_\mu=\partial_\mu\Lambda\) cumple \(F_{\mu\nu}=0\) y no ejerce ninguna fuerza electromagnética; entonces, una partícula libre se mueve a lo largo de una línea de universo recta. \lecturetimestamp{01:45:22}}
\end{classroomqa}

Por tanto, la invariancia gauge no es un mero truco prescindible. La electrodinámica, la teoría de Yang--Mills, el Modelo Estándar y la gravedad presentan redundancias gauge de un tipo u otro. El patrón es peculiar: se necesita una cantidad auxiliar para formular la teoría, pero parte de esa cantidad puede cambiarse sin alterar la física. Esa peculiaridad también es lo que permite organizarla.

\section{La lógica de la invariancia gauge}

La clase ha seguido un camino deliberadamente indirecto. Primero recuperamos el lenguaje de los índices arriba--abajo, porque las contracciones escalares son la gramática de las acciones invariantes de Lorentz. Después enunciamos cuatro principios: acción, localidad, invariancia de Lorentz e invariancia gauge. El cuarto quedó sin definir mientras reconstruíamos el cálculo de la partícula cargada.

El cálculo comenzó con la acción escalar local
\begin{equation}
I
=
-m\int d\tau
-e\int A_\mu dx^\mu,
\end{equation}
pasó por el lagrangiano en tiempo ordinario y por cada derivada de Euler--Lagrange, y recuperó la fuerza de Lorentz eléctrica y magnética. El tiempo propio reunió entonces el resultado en
\begin{equation}
m\frac{du^\mu}{d\tau}
=
-eF^\mu{}_\nu u^\nu,
\qquad
F_{\mu\nu}
=
\partial_\mu A_\nu-\partial_\nu A_\mu.
\end{equation}

Los ejemplos de campos escalares prepararon el último paso: una derivada total solo puede modificar la acción a través de los términos de contorno. En consecuencia,
\begin{equation}
A_\mu\longrightarrow A_\mu+\partial_\mu\Lambda
\end{equation}
solo cambia la acción de la línea de universo en los valores fijos de los extremos y deja \(F_{\mu\nu}\) inalterado porque las derivadas parciales mixtas conmutan. Por tanto, el potencial vectorial es necesario, pero no único. La invariancia gauge afirma que esa falta de unicidad pertenece a la descripción, no a la física.
